\chapter{运动：匀速与变速}\label{ch:g9:uniform-varied-motion}

两年前，你凭肉眼判断运动：点距相等，
即为匀速；点越散越开，即为加速。今年，点上
将加上毫米刻度，判断也要靠数字——逐段、
逐段计算，逐个速度比较。本章末尾有一句平静的话，
谈的是力与运动，日后将支撑起
物理学的半壁江山。

\section{运动的记录}

\begin{definition}[定时位置记录]\label{def:g9:uniform-varied-motion:record}
运动的\emph{定时位置记录}\index{定时位置记录}，是指
在等时间间隔处标记运动物体的位置——
视频逐帧读取、闪烁的频闪照片，
或实验室里打点计时器在被物体拉动的纸带上打下的点，
每秒 50 个点。
相邻两个标记之间，是一个时间间隔内的位移：
记录将运动化为可测量的线段。
\end{definition}

\begin{method}[用数字读取记录]\label{met:g9:uniform-varied-motion:read}
若记录的时间间隔为 $\tau$（例如每秒 50 点时取
\qty{0.02}{s}，视频中可取方便的 \qty{0.1}{s}）：
\begin{enumerate}
  \item 用毫米刻尺测量相邻两标记
        之间的间距；
  \item 将每段间距除以时间间隔：$v = d/\tau$ 给出
        该小段内的（平均）速度——间距以米计时，
        速度单位为 \unit{m/s}；
  \item 将各速度并排列出，读出其中的故事：
        数值平稳，为匀速运动；数值上升，为加速；
        数值下沉，为减速；
  \item 将任一小段的速度当作“该时刻”的速度：
        这段时间极短，平均速度与瞬时速度
        几乎相同。
\end{enumerate}
\end{method}

\begin{omfigure}
\begin{tikzpicture}[scale=0.85]
  \foreach \x in {0.5,2.0,3.5,5.0,6.5,8.0,9.5} {
    \fill[omDef] (\x,2.6) circle (0.1);
  }
  \node[right, font=\small] at (10.0,2.6)
    {间距均为 \qty{3.0}{cm}：匀速};
  \foreach \x in {0.5,1.1,2.0,3.2,4.7,6.5,8.6} {
    \fill[omThm] (\x,1.5) circle (0.1);
  }
  \node[right, font=\small] at (10.0,1.5)
    {间距渐增：加速};
  \foreach \x in {0.5,2.6,4.3,5.6,6.5,7.0,7.2} {
    \fill[omMeth] (\x,0.4) circle (0.1);
  }
  \node[right, font=\small] at (10.0,0.4)
    {间距渐减：减速};
\end{tikzpicture}

{\small 三条打点计时器纸带，每个时间间隔打一点。
今年测量各段间距，每段都化为一个
速度。}
\end{omfigure}

\begin{example}[纸带计算示例]\label{ex:g9:uniform-varied-motion:worked}
小车的纸带，时间间隔为 \qty{0.1}{s}，间距以
厘米计：依次为 $2.0$、$3.0$、$4.0$、$5.0$、$6.0$。逐段求速度：
$0.02 \div 0.1 = \qty{0.2}{m/s}$，随后为 $0.3$、$0.4$、
$0.5$、\qty{0.6}{m/s}。结论：加速——且
极为规律：速度每个时间间隔恰好增加
\qty{0.1}{m/s}。速度沿等量递增的运动，
即在相等时间内速度增量相等的运动，有专门的名称，
前景广阔——\emph{匀加速}运动，
是恒定原因的标志。
\end{example}

\begin{example}[最著名的加速运动]\label{ex:g9:uniform-varied-motion:freefall}
让石块下落并拍摄：间距逐个时间间隔拉大——
自由落体是加速运动。测得其速度每秒约增加
\qty{9.8}{m/s}，即下落每一秒增加约此值：一秒后，
速度为 \qty{9.8}{m/s}；两秒后，约为 \qty{20}{m/s}。
增长率等于当地的 $g$ 并非巧合——它
是今年物理中最深刻的呼应，高中分册
正是在此基础上构建其力学体系。记住它，
划出来，留待以后。
\end{example}

\section{瞬时速度的确切含义}

\begin{definition}[瞬时速度]\label{def:g9:uniform-varied-motion:instant}
运动在某一时刻的\emph{瞬时速度}\index{瞬时速度}，
是指在该时刻附近一段极短时间内的平均速度——
这段时间须短到
物体在其间来不及改变运动快慢。它就是
速度计所示的数值，也是雷达测得的数值：两者实际上都按
与\cref{met:g9:uniform-varied-motion:read}所用方法完全相同，
只是时间间隔极短，如同眨眼。（“极短”的确切含义，
须严格界定，是高中数学最后一年要解决的大问题——
这个答案催生了现代
科学。）
\end{definition}

\begin{example}[平均速度与瞬时速度对照]\label{ex:g9:uniform-varied-motion:sidebyside}
地铁两站之间的一次运行：平均速度——总
路程除以总时间——\qty{35}{km/h}。瞬时
速度：两站台处为 $0$，隧道中部为 \qty{70}{km/h}。
只有在匀速运动中，两者才
合为一个稳定的数值：匀速运动正是
使速度成为唯一、确切、不随时间间隔而变的量的运动。
\end{example}

\section{平静的一句话}

\begin{proposition}[平衡的力不改变运动]\label{prop:g9:uniform-varied-motion:inertia}
当物体所受的力相互抵消——或者根本
没有力作用——物体的运动不变：静止的
仍然静止；运动的则沿
直线以恒定速度继续运动。运动的改变——
加速、减速、转弯——只有在某个
\emph{非平衡力}作用时才会出现。此句先在此给出、
暂不证明，它是全部力学的开篇定律；高中
分册即从它开始，且始终以它为基础。
\end{proposition}

\begin{example}[用这句话读懂世界]\label{ex:g9:uniform-varied-motion:reading}
光滑冰面上的冰球笔直、匀速地滑行——受力
平衡（重力向下，冰面支持力向上），运动不变：无需推力
即可\emph{维持}滑行。刹车的汽车减速：有一个
非平衡力，即路面作用于轮胎的摩擦力，
方向向后。转弯的骑车人：侧向的非平衡力，
由倾斜的轮胎提供。而巡航中的客机，
以稳定的\qty{900}{km/h}飞行：推力平衡阻力，升力
平衡重力——四个力恰好相抵，匀速运动。
在物理学看来，静止与巡航
是同一种平静的状态。
\end{example}

\begin{remark}[这句话为何令人惊讶]\label{rem:g9:uniform-varied-motion:surprise}
日常经验似乎在反对：停止蹬踏，自行车
便会减速——运动难道不\emph{需要}力吗？但减速中的
自行车并非不受力：摩擦力与空气阻力向后推，
构成非平衡力，这正是速度衰减的原因。去掉
它们——比如冰场、太空真空——运动便
永远滑行下去：比你祖父母相识之前发射的
太空探测器至今仍在滑行，发动机早已熄火。
人类花了两千年才看穿摩擦力的
伪装；你们则有冰场与太空探测器
这样的优势。
\end{remark}

\section{练习}

\begin{exercise}[$\star$]\label{exo:g9:uniform-varied-motion:1}
什么是频闪记录？请说出
两种制作方法。
\end{exercise}

\begin{exercise}[$\star$]\label{exo:g9:uniform-varied-motion:2}
一条纸带，计时点间隔为 \qty{0.1}{s}，间距依次为 \qty{4.0}{cm}，
$4.0$，$4.0$，$4.0$。判断该运动，
并求其速度，单位为 \unit{m/s}。
\end{exercise}

\begin{exercise}[$\star$]\label{exo:g9:uniform-varied-motion:3}
另一条纸带：间距依次为 $1.0$，$2.5$，$4.5$，\qty{7.0}{cm}。
判断其运动性质，并计算首段与末段
的区间速度。
\end{exercise}

\begin{exercise}[$\star$]\label{exo:g9:uniform-varied-motion:4}
结合一次上学途中的情况，区分平均速度与瞬时
速度：校门口的测速雷达读出的是哪一个？家人所说的
“从家门到校门二十分钟”算出的是哪一个？
\end{exercise}

\begin{exercise}[$\star$]\label{exo:g9:uniform-varied-motion:5}
在哪些运动中，平均速度与瞬时速度
相等？为什么只有在那里才相等？
\end{exercise}

\begin{exercise}[$\star$]\label{exo:g9:uniform-varied-motion:6}
自由下落的石块，速度每秒增加多少，
这一增量等于当地的哪个数值？
\end{exercise}

\begin{exercise}[$\star$]\label{exo:g9:uniform-varied-motion:7}
请叙述那句平静的命题。它对一个
完全不受力的物体有何说法，以及对巡航客机
的四力平衡又有何说法？
\end{exercise}

\begin{exercise}[$\star$]\label{exo:g9:uniform-varied-motion:8}
“运动需要力才能持续。”请用滑行的
太空探测器与摩擦力的伪装，驳斥这一古老错误。
%
\end{exercise}

\begin{exercise}[$\star\star$]\label{exo:g9:uniform-varied-motion:9}
一条纸带（计时点间隔 \qty{0.1}{s}）记录的数据（单位：厘米）依次为 $6.0$，
$5.0$，$4.0$，$3.0$，$2.0$。判断其运动性质，计算每段
区间速度，并预测若此规律保持，纸带后续的数据。
\end{exercise}

\begin{exercise}[$\star\star$]\label{exo:g9:uniform-varied-motion:10}
例题中的小车，每个\qty{0.1}{s}的计时间隔，
速度增加 \qty{0.1}{m/s}。试用\emph{每秒}表示这一增量，并
将小车的增量与自由落体的增量相比较。
\end{exercise}

\begin{exercise}[$\star\star$]\label{exo:g9:uniform-varied-motion:11}
跳伞者越落越快——随后伞
张开，以稳定的 \qty{5}{m/s} 下落。用那句平静的命题
读懂两个阶段：开始时哪个力不平衡，最后
又形成了怎样的平衡？
\end{exercise}

\begin{exercise}[$\star\star\star$]\label{exo:g9:uniform-varied-motion:12}
请设计一份完整的实验结论：一辆玩具车从斜面
滑下并驶上地毯，应记录的数据、需计算的
数值、预期的两阶段判断（哪一阶段在
斜面上，哪一阶段在地毯上），以及用那句
关于力的命题解读每一阶段。
\end{exercise}

\section{问题：胶带局}

\begin{problem}[{周末问题 --- 运动分析局的一个下午；
四盘胶带，四个结论；
督察的顺口溜}]\label{pb:g9:uniform-varied-motion:1}
该局为学校、体育俱乐部和
法院分析计时记录。计时间隔始终为 \qty{0.1}{s}；间距以厘米给出。
你负责持尺。

\medskip\noindent\textbf{第一部分 --- 常规结论。}
\begin{enumerate}
  \item 胶带 A（走廊里的行人）：$8.0$，$8.0$，$8.0$，
        $8.0$，$8.0$。给出结论，并求速度，分别以 \unit{m/s}
        和 \unit{km/h} 表示。
  \item 胶带 B（短跑运动员的起跑）：$2.0$，$4.0$，$6.0$，
        $8.0$，$10.0$。给出结论，求首末速度，以及
        速度每秒的增量。
  \item 胶带 C（滚动的小球撞上沙子）：$10.0$，
        $7.0$，$4.5$，$2.5$，$1.0$。给出结论，并用
        关于力的命题说明沙子的作用。
  \item 胶带 D（一个谜）：$5.0$，$5.0$，$5.0$，$7.5$，
        $10.0$。讲述 D 的两阶段经过，并标出
        发生变化的那个计时点。
\end{enumerate}

\medskip\noindent\textbf{第二部分 --- 局里的
物理组。}
\begin{enumerate}[resume]
  \item 一位客户坚称其送货电动车“警官，平均速度是
        \qty{30}{km/h}”。雷达在路口记录到
        \qty{55}{km/h}。请向客户解释，
        为何两个数值都可以为真，以及
        法律以哪一个为准。
  \item 胶带 B 中的短跑运动员，速度以相等的量逐步增加。
        局里给这类运动贴上的是哪一个词，
        而这种增量的稳定性说明了推动力的
        什么特点？
  \item 胶带 A 中的走廊行人，据关于力的命题，是在平衡的力下行走的。
        请说出稳步行走者的平衡关系：
        什么在推，
        什么在阻。
  \item 某石块下落的胶带显示各段的平均速度为 $0.98$，
        $1.96$，\qty{2.94}{m/s}。请核对速度每秒的增量，
        并说出它再现的是哪一个数值。
\end{enumerate}

\medskip\noindent\textbf{第三部分 --- 法庭案件。}
一块滑板从斜坡滑下，穿过校园院子，
撞进花坛；管理员指责其“鲁莽的速度”，
滑板手则辩称“我一直在减速”。院子里的摄像头
记录了间距：斜坡出口为 $12.0$，随后依次为 $11.5$，$11.0$，
$10.5$，$10.0$，即进入花坛途中的间距。
\begin{enumerate}[resume]
  \item 对滑板手说法的结论，并给出
        各段的平均速度。
  \item 求出口速度与进入花坛时的速度，并以
        \unit{km/h} 表示：二者中是否有一个称得上“鲁莽”，
        与一位轻快骑车人的 \qty{15}{km/h} 相比如何？
  \item 管理员问，为何滑板不在平地上简单地
        停下，“事物本来就是如此”。
        局里的回答，请用
        一句关于力的语言给出。
  \item 以督察的顺口溜结束本案卷，两行：
        “间距相等意味着什么，
        间距改变又要求什么——依那句安静的命题，即一个原因。”
\end{enumerate}
\end{problem}
